\section{Robots and tasks}
\label{sec:robots-tasks}

\begin{lede}
What a robot remembers from one timestep to the next, and where the
never-ending stream of jobs comes from.
\end{lede}

\subsection{Robot state}

\code{Robot} carries, among other fields: \code{priority} (recomputed every
step, \Cref{sec:priority}); \code{tie\_breaker} $= 10^{-4} \cdot \mathrm{id}$,
deterministic, unique, and small enough never to reorder two robots across a
priority-class boundary; three related but distinct stall counters
\code{waiting\_time}, \code{blocked\_time} and \code{no\_progress\_steps}
(\Cref{sec:deadlock-signals} distinguishes them precisely, because they feed
different formulas); \code{route}, the current planned path;
\code{route\_distance\_to\_waypoint} $= \Dg{w}(\pos{i}{t})$, the live BFS
distance to the current waypoint; \code{mode}, a proximity bucket
(\Cref{sec:score-mode}); \code{direction\_weight} and \code{aisle\_weight}
($\beta_i$, $\gamma_i$, \Cref{sec:score-beta,sec:score-gamma}); and
\code{position\_history}, a 16-entry deque used for repeated-configuration
detection (\Cref{sec:deadlock-detection}).

\subsection{Task urgency and waiting}

\begin{align}
  \mathrm{urgency}(t) &=
    \begin{cases}
      0 & \text{if } \mathrm{deadline} = \INF,\\[2pt]
      \dfrac{1}{\max(1,\ \mathrm{deadline} - t)} & \text{otherwise,}
    \end{cases}\\[4pt]
  \mathrm{waiting\_time}(t) &= \max(0,\ t - t_{\mathrm{rel}}),\\
  \mathrm{service\_time} &= t_{\mathrm{complete}} - t_{\mathrm{rel}}
    \quad \text{(defined only once completed).}
\end{align}

Urgency grows without bound as the deadline approaches and passes
(slack $\rightarrow 0$ or negative), which is intentional: a task past its
deadline should dominate the priority function's urgency term
(\Cref{sec:priority}) rather than saturate.

\subsection{Arrival processes}
\label{sec:arrivals}

Four modes, each returning a task count for \code{receive\_new\_tasks(t)}:
\begin{align*}
  \text{\code{periodic}:}\ & \text{count} = \mathrm{rate} \ \text{ if } t \bmod \mathrm{period} = 0,\ \text{else } 0;\\
  \text{\code{poisson}:}\ & \text{count} \sim \mathrm{Poisson}(\lambda) \text{ via Knuth's sampler};\\
  \text{\code{bursty}:}\ & \text{count} = \mathrm{burst\_size} \ \text{ if } t \bmod \mathrm{burst\_period} = 0,\ \text{else } 0;\\
  \text{\code{batch}:}\ & \text{count} = \mathrm{total} \ \text{ if } t = 0,\ \text{else } 0.
\end{align*}

Knuth's Poisson sampler is the standard method for small-to-moderate $\lambda$:

\begin{algorithm}[H]
\caption{Knuth's Poisson sampler (\code{task.\_poisson})}
\begin{algorithmic}[1]
\Require mean $\lambda > 0$
\State $L \gets e^{-\lambda}$;\quad $k \gets 0$;\quad $p \gets 1$
\Loop
  \State $p \gets p \cdot U$ \Comment{$U \sim \mathrm{Uniform}(0,1)$}
  \If{$p \leq L$} \Return $k$ \EndIf
  \State $k \gets k + 1$
\EndLoop
\end{algorithmic}
\end{algorithm}

After $k$ multiplications $p = \prod_{1}^{k} U_j$, and the number of uniform
draws needed to push that running product below $e^{-\lambda}$ is exactly
Poisson distributed with mean $\lambda$ \citep{knuth1997taocp2}. The
implementation guards at $k > 1000$ for numerical safety, which never triggers
at realistic rates.
